\twocolumn

\section{Problem Statement}
\label{sec:problem}

Our goal is to present a deterministic password generator (DPG) design that remedies the security, privacy, and usability issues of existing DPGs. In this section, we detail the specific goals and assumptions of the presented MFDPG construction.

\subsection{Threat Model}
Consider a DPG implemented as a client-side application or browser extension with optional server-side storage of settings, salts, and other trustless public parameters. 
We consider security under a \textit{total data breach} threat model; i.e., at some instant, an adversary receives a snapshot of all materials stored on the client or server. The adversary may use all such data to attempt to violate one or more of the below security goals.

% Thus, our design must first and foremost satisfy the following standard properties of DPGs. Though existing DPGs use a master password as a sole input factor, we generalize these definitions to support a collection of input factors for forward-compatibility with our scheme. Let one such set of input factors be defined as ``correct'' for a given user. The desired properties are then as follows:

% data breach of any stored parameters
% adversaries have access to generated passwords

\subsection{Security Goals}
\label{sec:goals}

Our design must first and foremost satisfy the following standard properties of DPGs. Though existing DPGs use a master password as a sole input factor, we generalize these definitions to support a collection of input factors for forward-compatibility with our scheme. Let one such set of input factors (and corresponding output) be defined as ``correct'' for a given user. The desired properties are then as follows:

\begin{enumerate}[leftmargin=*]
    \itemsep 1em
    \item \textbf{Correctness} (Determinism). Given a static set of ``correct'' input factors, the scheme always outputs the same ``correct'' password for any target service; i.e., it is deterministic. 
    \item \textbf{Safety} (Pseudorandomness). Given an ``incorrect'' set of input factors, the scheme outputs an ``incorrect'' password for any target service except with negligible probability. 
    \item \textbf{Secretless}. While password managers can achieve the above properties through stateful storage of ciphertexts, a DPG should not store site-specific ciphertexts in any location, i.e., it is stateless other than public parameters.
\end{enumerate}

Because the goal of this paper is to significantly improve upon the current state-of-the-art DPGs, we impose a few additional requirements on our MFDPG scheme:

\begin{enumerate}[resume,leftmargin=*]
    \itemsep 1em
    \item \textbf{Brute-Force Resistance}. A brute-force attack on the master password should be insufficient to generate site-specific passwords. Adversaries that obtain a site-specific password should not easily be able to attack the master password.
    \item \textbf{Compatibility}. The DPG should be compatible with a wide variety of password policies. Namely, the DPG should be able to generate passwords compliant with any regular password policy (i.e., representable by a regular expression).
    \item \textbf{Revocability}. Users should, at any time, be able to revoke a password for a given service and generate a new one for the same service, without changing their input factors, and without manually remembering a counter value.
    \item \textbf{Privacy}. In implementing the above properties, a DPG should not store site-specific values that allow any outside adversary to identify the services utilized by the DPG user.
\end{enumerate}

\eject